
\documentclass{article}
\usepackage{colortbl}
\usepackage{makecell}
\usepackage{multirow}
\usepackage{supertabular}

\begin{document}

\newcounter{utterance}

\centering \large Interaction Transcript for game `referencegame', experiment `line\_grids\_rows', episode 3 with NEW\_Reflect\_Agent.
\vspace{24pt}

{ \footnotesize  \setcounter{utterance}{1}
\setlength{\tabcolsep}{0pt}
\begin{supertabular}{c@{$\;$}|p{.15\linewidth}@{}p{.15\linewidth}p{.15\linewidth}p{.15\linewidth}p{.15\linewidth}p{.15\linewidth}}
    \# & $\;$A & \multicolumn{4}{c}{Game Master} & $\;\:$B\\
    \hline

    \theutterance \stepcounter{utterance}  
    & & \multicolumn{4}{p{0.6\linewidth}}{
        \cellcolor[rgb]{0.9,0.9,0.9}{
            \makecell[{{p{\linewidth}}}]{
                \texttt{\tiny{[P1$\langle$GM]}}
                \texttt{You are given three grids, where each of them is 5 by 5 in size.} \\
\texttt{Grids have empty cells marked with "▢" and filled cells marked with "X".} \\
\texttt{Your task is to generate a referring expression that best describes the target grid while distinguishing it from the two other distractor grids.} \\
\texttt{The first grid is the target grid, and the following two grids are the distractors.} \\
\\ 
\texttt{Target grid:} \\
\\ 
\texttt{□ □ X X X} \\
\texttt{□ □ X X X} \\
\texttt{□ □ X X X} \\
\texttt{□ □ X X X} \\
\texttt{□ □ X X X} \\
\\ 
\texttt{Distractor grid 1:} \\
\\ 
\texttt{□ □ □ X X} \\
\texttt{□ □ □ X X} \\
\texttt{□ □ □ X X} \\
\texttt{□ □ □ X X} \\
\texttt{□ □ □ X X} \\
\\ 
\texttt{Distractor grid 2:} \\
\\ 
\texttt{X □ □ X X} \\
\texttt{X □ □ X X} \\
\texttt{X □ □ X X} \\
\texttt{X □ □ X X} \\
\texttt{X □ □ X X} \\
\\ 
\texttt{Instruction: Describe the target grid.} \\
\texttt{Generate the referring expression starting with the tag "Expression: " for the given target grid. Omit any other text.} \\
            }
        }
    }
    & \\ \\

    \theutterance \stepcounter{utterance}  
    & \multicolumn{4}{p{0.45\linewidth}}{
        \cellcolor[rgb]{0.95,0.95,0.95}{
            \makecell[{{p{\linewidth}}}]{
                \texttt{\tiny{[P1$|$internal]}}
                \texttt{Reflections from previous episodes:} \\
\texttt{{-}} \\
\texttt{ Expression: The grid has two full rows of X's at the top and bottom, with all other rows empty.} \\
\\ 
\texttt{Reflection: Failed. The description did not distinguish the target from the distractors, as both distractors also have rows with all X's. I should have specified the exact positions (first and fourth rows) and emphasized that only those two rows are filled, with the rest completely empty, unlike the distractors which have vertical blocks. Next time, I will explicitly note the row numbers and contrast the horizontal full rows versus vertical partial columns.} \\
\texttt{{-}} \\
\texttt{ Failed. My expression "two full rows of X's at the top and bottom" was ambiguous—it didn't distinguish from distractors, which also have full rows of X's. I focused on the target's pattern but neglected a precise comparison. Next time, I must explicitly note the unique feature: the target is the only grid with **two separate, complete horizontal rows of X's** (rows 1 and 4) with no other X's.} \\
\texttt{{-}} \\
\texttt{ Failed. I described the target as having "two full rows of Xs" but didn't specify they were the first and fourth rows. This was insufficient to distinguish it from the distractors, which also had rows with multiple Xs. Next time, I must specify the exact positions (rows 1 and 4) and note the completely empty rows in between.} \\
            }
        }
    }
    & & \\ \\

    \theutterance \stepcounter{utterance}  
    & \multicolumn{4}{p{0.45\linewidth}}{
        \cellcolor[rgb]{0.95,0.95,0.95}{
            \makecell[{{p{\linewidth}}}]{
                \texttt{\tiny{[P1$|$internal]}}
                \texttt{{[}attempt 1 – valid{]}} \\
\texttt{Expression: The grid has two full rows of X's at the top and bottom, with all other rows empty.} \\
            }
        }
    }
    & & \\ \\

    \theutterance \stepcounter{utterance}  
    & \multicolumn{4}{p{0.6\linewidth}}{
        \cellcolor[rgb]{0.9,0.9,0.9}{
            \makecell[{{p{\linewidth}}}]{
                \texttt{\tiny{[P1$\rangle$GM]}}
                \texttt{Expression: The grid has two full rows of X's at the top and bottom, with all other rows empty.} \\
            }
        }
    }
    & & \\ \\

    \theutterance \stepcounter{utterance}  
    & & & \multicolumn{2}{p{0.3\linewidth}}{
        \cellcolor[rgb]{0.9,0.9,0.9}{
            \makecell[{{p{\linewidth}}}]{
                \texttt{\tiny{[GM$|$GM]}}
                \texttt{Expression: The grid has two full rows of X's at the top and bottom, with all other rows empty.} \\
            }
        }
    }
    & & \\ \\

    \theutterance \stepcounter{utterance}  
    & & \multicolumn{4}{p{0.6\linewidth}}{
        \cellcolor[rgb]{0.9,0.9,0.9}{
            \makecell[{{p{\linewidth}}}]{
                \texttt{\tiny{[GM$\rangle$P2]}}
                \texttt{You are given three grids, where each of them is 5 by 5 in size.} \\
\texttt{Grids have empty cells marked with "▢" and filled cells marked with "X".} \\
\texttt{You are also given a referring expression that describes one of the given grids.} \\
\texttt{Your task is to select the grid that matches the given referring expression.} \\
\texttt{Generate only the number (in text) of the grid that matches the given expression by selecting first, second, or third.} \\
\\ 
\texttt{First:} \\
\\ 
\texttt{□ □ X X X} \\
\texttt{□ □ X X X} \\
\texttt{□ □ X X X} \\
\texttt{□ □ X X X} \\
\texttt{□ □ X X X} \\
\\ 
\texttt{Second:} \\
\\ 
\texttt{□ □ □ X X} \\
\texttt{□ □ □ X X} \\
\texttt{□ □ □ X X} \\
\texttt{□ □ □ X X} \\
\texttt{□ □ □ X X} \\
\\ 
\texttt{Third:} \\
\\ 
\texttt{X □ □ X X} \\
\texttt{X □ □ X X} \\
\texttt{X □ □ X X} \\
\texttt{X □ □ X X} \\
\texttt{X □ □ X X} \\
\\ 
\texttt{Expression: The grid has two full rows of X's at the top and bottom, with all other rows empty.} \\
\texttt{Question: Which grid does the expression refer to?} \\
\texttt{Start with the tag "Answer: ", followed by your selection. Omit any other text.} \\
            }
        }
    }
    & \\ \\

    \theutterance \stepcounter{utterance}  
    & & & \multicolumn{4}{p{0.45\linewidth}}{
        \cellcolor[rgb]{0.95,0.95,0.95}{
            \makecell[{{p{\linewidth}}}]{
                \texttt{\tiny{[P2$|$internal]}}
                \texttt{Reflections from previous episodes:} \\
\texttt{{-}} \\
\texttt{ Failed. I incorrectly selected the second grid, but the expression specifies a solid top row and a solid *fourth* row, with all other rows empty. The second grid has a solid top row and a solid *fifth* row. I misread the row numbering. Next time, I will count rows precisely from 1 to 5 and verify each condition against the correct row index.} \\
\texttt{{-}} \\
\texttt{ Answer: First} \\
\\ 
\texttt{Reflection: Succeeded. I correctly matched the expression to Grid 1 because it has two full rows of X's (first and fourth rows) and all other rows empty, with no full columns. I verified each grid against the description. Next time, I will continue to carefully parse each condition in the expression and check all grids systematically.} \\
\texttt{{-}} \\
\texttt{ Answer: Third} \\
\\ 
\texttt{Reflection: Succeeded. Correctly matched the expression's description of two full horizontal rows of X's in rows 1 and 4, with all other cells empty. The analysis focused on verifying each grid against the specific row conditions, ignoring irrelevant column details. Next time, maintain the same method of checking row patterns precisely against the expression.} \\
            }
        }
    }
     \\ \\

    \theutterance \stepcounter{utterance}  
    & & & \multicolumn{4}{p{0.45\linewidth}}{
        \cellcolor[rgb]{0.95,0.95,0.95}{
            \makecell[{{p{\linewidth}}}]{
                \texttt{\tiny{[P2$|$internal]}}
                \texttt{{[}attempt 1 – valid{]}} \\
\texttt{Answer: First} \\
            }
        }
    }
     \\ \\

    \theutterance \stepcounter{utterance}  
    & & & \multicolumn{4}{p{0.6\linewidth}}{
        \cellcolor[rgb]{0.9,0.9,0.9}{
            \makecell[{{p{\linewidth}}}]{
                \texttt{\tiny{[GM$\langle$P2]}}
                \texttt{Answer: First} \\
            }
        }
    }
     \\ \\

    \theutterance \stepcounter{utterance}  
    & & & \multicolumn{2}{p{0.3\linewidth}}{
        \cellcolor[rgb]{0.9,0.9,0.9}{
            \makecell[{{p{\linewidth}}}]{
                \texttt{\tiny{[GM$|$GM]}}
                \texttt{first} \\
            }
        }
    }
    & & \\ \\

\end{supertabular}
}

\end{document}
